\documentclass[conference]{IEEEtran}
\IEEEoverridecommandlockouts

\usepackage{cite}
\usepackage{amsmath,amssymb}
\usepackage{graphicx}
\usepackage{booktabs}
\usepackage[brazil]{babel}
\usepackage[utf8]{inputenc}
\usepackage[T1]{fontenc}
\usepackage{url}
\usepackage{hyperref}

\graphicspath{{figuras/}}

\begin{document}

\title{VAE ou GAN para aumento de dados em radiografias de tórax?\\
Um estudo pareado com 30 simulações}

\author{
\IEEEauthorblockN{Guilherme Alencar}
\IEEEauthorblockA{\textit{CESAR School} \\
Recife, Brasil}
\and
\IEEEauthorblockN{Leonardo Mello}
\IEEEauthorblockA{\textit{CESAR School} \\
Recife, Brasil}
}

\maketitle

\begin{abstract}
O aumento de dados com modelos generativos é frequentemente apresentado como
solução para classes minoritárias escassas em imagens médicas, mas raramente
se compara geradores sob controle experimental rigoroso. Este trabalho compara
duas abordagens que diferem em uma única técnica: um Autoencoder Variacional
(VAE) e uma DCGAN, ambos usados para sintetizar radiografias da classe
minoritária. O restante do pipeline é idêntico nos dois cenários: ResNet18
pré-treinada e congelada como extratora de características, seguida de um
ensemble de votação suave. Avaliamos no PneumoniaMNIST com 30 simulações
pareadas, em que cada repetição compartilha semente, dados reais e conjunto de
teste. Rejeitamos a hipótese nula: o cenário com DCGAN obteve F1 da classe
minoritária superior ao do VAE ($\Delta = 0{,}032$; $p < 10^{-6}$;
$d_z = 2{,}40$). O tamanho de efeito é grande, mas o ganho absoluto é modesto e
o cenário com DCGAN apenas empata com a linha de base sem aumento algum
($0{,}717$ contra $0{,}710$), enquanto o VAE a degrada. Uma análise de
separabilidade no espaço de características mostra por quê: um classificador
linear distingue amostras sintéticas de reais com acurácia superior a
$0{,}98$ para ambos os geradores, de modo que o balanceamento obtido é nominal
e não efetivo. Concluímos que significância estatística e relevância prática
divergem neste problema, e que verificar a separabilidade das amostras
sintéticas deveria preceder qualquer alegação de ganho por aumento generativo.
\end{abstract}

\begin{IEEEkeywords}
aumento de dados, modelos generativos, VAE, GAN, transfer learning,
aprendizado por ensemble, classes desbalanceadas, imagem médica
\end{IEEEkeywords}

\section{Introdução}

Classificadores de imagens médicas enfrentam, com frequência, uma distribuição
de classes desbalanceada: achados raros são, por definição, pouco
representados. Em radiografias de tórax para triagem de pneumonia, um serviço
que inicia a rotulação de seu acervo costuma dispor de poucos exames normais
confirmados frente a um volume maior de casos positivos. Um classificador
treinado nesse regime tende a favorecer a classe majoritária, e a métrica de
acurácia mascara o problema.

O aumento de dados com modelos generativos é uma resposta recorrente:
sintetizam-se exemplos da classe escassa até equilibrar o treino
\cite{frid2018gan}. A literatura, porém, costuma reportar o ganho de um único
gerador contra uma linha de base, sem comparar geradores entre si sob o mesmo
protocolo, e sem verificar se as amostras sintéticas de fato se integram à
distribuição real.

Este trabalho trata exatamente dessa lacuna. Comparamos dois cenários que
diferem em \emph{uma única} técnica, mantendo todo o resto constante:

\begin{itemize}
  \item \textbf{Cenário H0:} VAE convolucional \cite{kingma2014vae} como gerador;
  \item \textbf{Cenário H1:} DCGAN \cite{radford2016dcgan} como gerador.
\end{itemize}

Ambos alimentam o mesmo classificador: uma ResNet18 \cite{he2016resnet}
pré-treinada no ImageNet e congelada, seguida de um ensemble de votação suave
\cite{dietterich2000ensemble}. As contribuições são:

\begin{enumerate}
  \item um protocolo pareado de 30 simulações que isola o efeito do gerador,
        com teste de hipótese, tamanho de efeito e intervalo de confiança;
  \item a constatação de que a diferença entre geradores, embora
        estatisticamente muito significativa, é praticamente modesta frente à
        linha de base sem aumento;
  \item um diagnóstico de separabilidade que explica o resultado e que
        propomos como verificação mínima antes de alegar ganho por aumento
        generativo.
\end{enumerate}

\section{Trabalhos relacionados}

O desbalanceamento de classes é tradicionalmente tratado por reamostragem,
sendo o SMOTE \cite{chawla2002smote} o representante clássico ao interpolar
vizinhos no espaço de atributos. Com imagens, a interpolação linear em pixels
produz artefatos, o que motivou o uso de modelos generativos profundos.

VAEs \cite{kingma2014vae} otimizam um limite inferior da verossimilhança e
produzem um espaço latente contínuo, do qual se amostra por construção. Por
usarem perda de reconstrução pixel a pixel, tendem a gerar imagens suavizadas.
GANs \cite{goodfellow2014gan} substituem a reconstrução por um discriminador
adversarial; a DCGAN \cite{radford2016dcgan} estabeleceu as escolhas
arquiteturais que tornaram esse treino estável. Frid-Adar et
al.~\cite{frid2018gan} relatam ganho com síntese por GAN em lesões hepáticas,
um dos resultados que motivam a prática.

Para avaliar fidelidade de geradores, a distância de Fréchet entre ativações de
uma rede pré-treinada \cite{heusel2017fid} tornou-se padrão. Neste trabalho a
calculamos sobre as características da própria ResNet18 do pipeline, e a
reportamos ao lado do desempenho de classificação para examinar se as duas
medidas concordam.

\section{Metodologia}

\subsection{Conjunto de dados}

Utilizamos o PneumoniaMNIST, do benchmark MedMNIST v2 \cite{yang2023medmnist}:
radiografias de tórax pediátricas em $28 \times 28$, tons de cinza, com rótulos
\emph{normal} e \emph{pneumonia}. A Tabela~\ref{tab:dados} descreve a partição
oficial, preservada integralmente.

\begin{table}[t]
\centering
\caption{Partição oficial do PneumoniaMNIST}
\label{tab:dados}
\begin{tabular}{lccc}
\toprule
Partição & Imagens & Normal & Pneumonia \\
\midrule
Treino & 4.708 & 1.214 (25,8\%) & 3.494 (74,2\%) \\
Teste  & 624   & 234 (37,5\%)   & 390 (62,5\%)   \\
\bottomrule
\end{tabular}
\end{table}

\subsection{Escassez simulada}

Com as 1.214 imagens normais disponíveis, o aumento sintético teria pouco a
acrescentar. Para construir o regime em que a geração é de fato necessária,
cada simulação sorteia, sem reposição, $300$ imagens normais e $1.500$ com
pneumonia, uma razão de $1$:$5$. Cada gerador sintetiza $1.200$ imagens
normais, levando o treino a $1.500$ por classe.

\subsection{Pipeline}

A Fig.~\ref{fig:pipeline} resume o fluxo. As três técnicas são:

\textbf{A. Geração sintética.} VAE convolucional (cenário H0) ou DCGAN
(cenário H1). Esta é a \emph{única} diferença entre os cenários.

\textbf{B. Transfer learning.} ResNet18 pré-treinada no ImageNet, com a camada
de classificação substituída pela identidade, produzindo um vetor de $512$
características. A rede permanece \emph{congelada} por três razões: com $300$
imagens reais, ajustar $11$ milhões de parâmetros levaria a sobreajuste; o
custo de $90$ ajustes finos seria proibitivo; e, sobretudo, manter a extratora
fixa garante que a única fonte de variação entre cenários seja o gerador. As
imagens são reamostradas para $112 \times 112$ e replicadas em três canais,
aproximando a escala e o formato esperados pela rede.

\textbf{C. Ensemble.} Votação suave entre regressão logística, \emph{random
forest} e \emph{gradient boosting} histogramado, sobre as características
extraídas. Os três modelos usam ponderação de classes inversamente
proporcional à frequência, aplicada \emph{identicamente} nos três cenários;
trata-se de correção do instrumento de medida, não do tratamento.

\begin{figure}[t]
\centering
% Substituir por diagrama do pipeline (ver paper/README.md)
\framebox[0.95\columnwidth]{\rule{0pt}{4.2cm}\ \ Diagrama do pipeline}
\caption{Fluxo experimental. Apenas o bloco de geração sintética difere entre
os cenários H0 e H1; extração de características, classificador e conjunto de
teste são idênticos.}
\label{fig:pipeline}
\end{figure}

\subsection{Protocolo experimental}

São $30$ simulações independentes, com sementes $42$ a $71$. A semente controla
o sorteio das imagens reais, a inicialização dos geradores, a ordem dos lotes e
a aleatoriedade do ensemble. Como os dois cenários compartilham a mesma
semente, cada par observa \emph{as mesmas} imagens reais e é avaliado no mesmo
conjunto de teste oficial, o que caracteriza um desenho pareado. Para que
nenhum gerador seja favorecido por mais computação, ambos recebem orçamento
idêntico: $200$ épocas. Um terceiro cenário, \emph{sem aumento}, é registrado
como linha de base e não participa do teste.

\subsection{Métricas e análise estatística}

A métrica principal, definida antes da execução, é o F1 da classe minoritária.
A acurácia seria enganosa: o teste contém $62{,}5\%$ de pneumonia, de modo que
um classificador trivial já a atingiria sem aprender nada. Reportamos ainda
AUC, acurácia balanceada, revocação e a distância de Fréchet sobre as
características da ResNet18 (FD), esta última como medida de fidelidade
independente do classificador.

As hipóteses sobre a diferença média no F1 são
$H_0: \mu_{\text{VAE}} - \mu_{\text{GAN}} = 0$ contra
$H_1: \mu_{\text{VAE}} - \mu_{\text{GAN}} \neq 0$. O teste é bilateral, pois
não havia expectativa fundamentada sobre qual gerador venceria. Verificamos a
normalidade das $30$ diferenças pareadas com Shapiro--Wilk; satisfeita a
suposição, aplicamos o teste $t$ pareado, e o de Wilcoxon caso contrário.
Adotamos $\alpha = 0{,}05$ e reportamos o tamanho de efeito $d_z$ de Cohen e o
intervalo de confiança de $95\%$.

\section{Resultados}

\subsection{Desempenho de classificação}

A Tabela~\ref{tab:resultados} apresenta média e desvio padrão nas $30$
simulações. O cenário com DCGAN supera o com VAE no F1 da classe minoritária, e
o teste pareado rejeita $H_0$ com folga: diferença média de $-0{,}032$ a favor
da DCGAN, $p < 10^{-6}$, $d_z = -2{,}40$, intervalo de confiança de $95\%$
$[-0{,}037;\ -0{,}027]$. Shapiro--Wilk não rejeitou a normalidade das
diferenças ($p = 0{,}69$), de modo que o teste $t$ pareado é apropriado.

% Tabela gerada automaticamente em resultados/tabela_resultados.tex
\begin{table}[t]
\centering
\caption{Média (desvio padrão) em 30 simulações pareadas}
\label{tab:resultados}
\begin{tabular}{lccc}
\toprule
Métrica & Sem aumento & VAE (H0) & GAN (H1) \\
\midrule
F1 (minoritária)      & 0,710 & 0,685 & \textbf{0,717} \\
AUC                   & \textbf{0,939} & 0,928 & 0,919 \\
Acurácia balanceada   & 0,775 & 0,759 & \textbf{0,778} \\
Revocação (normal)    & 0,570 & 0,534 & \textbf{0,586} \\
FD ResNet ($\downarrow$) & --- & 223,3 & \textbf{160,2} \\
\bottomrule
\end{tabular}
\end{table}

Dois aspectos merecem destaque. Primeiro, a linha de base \emph{sem aumento
algum} alcança F1 de $0{,}710$, de modo que a DCGAN a supera por apenas
$0{,}007$, enquanto o VAE a degrada em $0{,}025$. Segundo, a AUC favorece a
linha de base sobre ambos os geradores, indicando que o aumento sintético
altera o limiar de decisão mais do que a capacidade de ordenação do modelo.

\subsection{Fidelidade das amostras}

A distância de Fréchet favorece a DCGAN ($160{,}2$ contra $223{,}3$), em
concordância com o desempenho de classificação. Neste experimento, portanto,
fidelidade e utilidade apontam na mesma direção.

\subsection{Diagnóstico de separabilidade}

Para verificar se o balanceamento numérico se traduz em balanceamento efetivo,
treinamos uma regressão logística para distinguir características de imagens
reais das de sintéticas, com validação cruzada. Se as amostras sintéticas
estivessem integradas à distribuição real, essa acurácia ficaria próxima de
$0{,}5$. A Tabela~\ref{tab:separabilidade} mostra o contrário.

\begin{table}[t]
\centering
\caption{Separabilidade entre características reais e sintéticas (semente 42)}
\label{tab:separabilidade}
\begin{tabular}{lccc}
\toprule
Gerador & Separabilidade & Norma média & \% da real \\
\midrule
Reais & --- & 34,4 & 100,0\% \\
VAE   & 0,997 & 27,3 & 79,4\% \\
GAN   & 0,986 & 34,6 & 100,5\% \\
\bottomrule
\end{tabular}
\end{table}

Ambos os geradores produzem amostras quase perfeitamente separáveis das reais
por uma fronteira \emph{linear}. Do ponto de vista do classificador, as imagens
sintéticas não constituem exemplos adicionais da classe minoritária, mas um
agrupamento próprio, que os modelos baseados em árvores isolam com poucas
partições. O balanceamento de $1$:$1$ é, portanto, nominal e não efetivo.

A norma média das características esclarece a diferença entre os geradores: a
DCGAN reproduz a magnitude de ativação das imagens reais ($100{,}5\%$),
enquanto o VAE fica substancialmente abaixo ($79{,}4\%$), compatível com a
suavização característica de sua perda de reconstrução. Isso explica por que a
DCGAN, apesar de igualmente separável, degrada menos o classificador.

\section{Discussão}

\subsection{Significância estatística e relevância prática}

O resultado ilustra com clareza a distinção entre as duas noções. O tamanho de
efeito $d_z = 2{,}40$ é, pelos critérios convencionais, muito grande, e o valor
$p$ é inferior a $10^{-6}$. Isso decorre do desenho pareado: ao comparar
geradores sobre as mesmas imagens, eliminamos a variância do sorteio e tornamos
detectável uma diferença sistemática pequena.

Pequena é a palavra adequada. Em termos absolutos, a troca do VAE pela DCGAN
rende $0{,}032$ de F1, e a DCGAN supera a ausência total de aumento por
$0{,}007$. Um serviço clínico que investisse em treinar um modelo generativo
obteria, neste cenário, ganho praticamente nulo. Reportar apenas o valor $p$
sugeriria uma conclusão bem mais entusiasmada do que os dados sustentam.

\subsection{Por que o aumento não entrega o ganho esperado}

O diagnóstico de separabilidade oferece a explicação. Treinar um gerador com
$300$ imagens de $28 \times 28$ não basta para reproduzir a estatística de
textura que a ResNet18 utiliza para descrever as radiografias reais. As
amostras sintéticas ocupam região própria do espaço de características, e um
classificador baseado em árvores simplesmente as isola, em vez de generalizar a
partir delas.

Essa constatação sugere uma prática metodológica: antes de atribuir ganho ao
aumento generativo, convém medir a separabilidade entre características reais e
sintéticas. Trata-se de verificação barata, que requer apenas um classificador
linear, e que distingue balanceamento nominal de balanceamento efetivo.

\subsection{Limitações}

A resolução de $28 \times 28$ limita o teto de qualidade dos geradores;
resoluções maiores plausivelmente reduziriam a separabilidade. A extratora
congelada favorece o controle experimental, mas o ajuste fino poderia adaptar
as características ao domínio radiográfico. O estudo cobre duas classes e um
único conjunto de dados, de modo que a generalização a outros domínios médicos
permanece em aberto. Por fim, parte das operações no backend MPS não é
determinística, o que introduz variação em casas decimais distantes, sem efeito
sobre as conclusões.

\subsection{Trabalhos futuros}

Três direções decorrem diretamente dos resultados: avaliar se resoluções
maiores reduzem a separabilidade a ponto de tornar o aumento efetivo;
contrastar a síntese generativa com aumento clássico por transformações
geométricas, de custo muito inferior; e investigar se filtrar amostras
sintéticas por um critério de proximidade às reais converte parte da síntese em
ganho mensurável.

\section{Conclusão}

Comparamos VAE e DCGAN como geradores para aumento da classe minoritária em
radiografias de tórax, sob protocolo pareado de $30$ simulações em que apenas o
gerador varia. Rejeitamos a hipótese nula: a DCGAN supera o VAE no F1 da classe
minoritária, com $p < 10^{-6}$ e tamanho de efeito grande, e também apresenta
maior fidelidade pela distância de Fréchet.

A leitura cuidadosa do resultado, contudo, modera o entusiasmo. A DCGAN apenas
empata com a ausência de aumento, e o VAE a degrada. O diagnóstico de
separabilidade explica o fenômeno: ambos os geradores produzem amostras quase
perfeitamente distinguíveis das reais, de modo que o balanceamento alcançado é
nominal. A contribuição prática deste trabalho é menos a escolha entre
geradores e mais a recomendação de verificar a separabilidade antes de atribuir
ganho ao aumento generativo.

\section*{Disponibilidade}

Código, documentação e instruções de reprodução estão disponíveis em
\url{https://github.com/GuilhermeAlencar22/vae-vs-gan-medical-augmentation}.

\bibliographystyle{IEEEtran}
\bibliography{referencias}

\end{document}
